\documentclass[11pt,compress,t,notes=noshow, xcolor=table]{beamer}

\input{../../style/preamble}
\input{../../latex-math/basic-math}
\input{../../latex-math/basic-ml}

\newcommand{\vb}{\mathbf{v}}
% centered footnotesize caption below a figure
\newcommand{\figcaption}[1]{\par{\centering\footnotesize #1\par}}

\title{Optimization in Machine Learning}

\begin{document}

\titlemeta{}{
Nelder-Mead method
}{
figure/nelder_mead_steps_3.png
}{
\item General idea
\item Reflection, expansion, contraction
\item Advantages \& disadvantages
\item Examples
}

\begin{framei}{Nelder-Mead method}
\item Derivative-free method $\Rightarrow$ heuristic
\item Generalization of bisection in $d$-dimensional space
\item Based on $d$-simplex, defined by $d + 1$ points:
\begin{itemize}
\item $d = 1$ interval
\comment[NOTE]{BB}{bilder malen}
\item $d = 2$ triangle
\item $d = 3$ tetrahedron
\item $\cdots$
\end{itemize}
\end{framei}


\begin{framev}{nelder-mead method}
A version of the Nelder-Mead method:\\
\spacer
Initialization: choose $d + 1$ random, affinely independent points $\vb_i$\\ ($\vb_i$ are vertices: corner points of the simplex/polytope)
\spacer
\begin{enumerate}
\item[1.] Order points according to ascending function values
$$f(\vb_1) \leq f(\vb_2) \leq \ldots \leq f(\vb_d) \leq f(\vb_{d + 1})$$
with $\vb_1$ best point, $\vb_{d + 1}$ worst point
\end{enumerate}
\imageC[0.34]{figure/nelder_mead_steps_1.png}
\figcaption{Example: Nelder-Mead in $d = 2$ dimensions with $d + 1 = 3$ initialized points}
\end{framev}


\begin{framev}{nelder-mead method}
\begin{enumerate}
\item[2.] Compute centroid without worst point
$$\bar{\vb} = \frac{1}{d} \sum_{i = 1}^d \vb_i.$$
\end{enumerate}
\imageC[0.42]{figure/nelder_mead_steps_2.png}
\figcaption{Simplex with computed centroid $\bar{\vb}$}
\end{framev}


\begin{framev}{nelder-mead method}
\begin{enumerate}
\item[3.] Reflection: compute reflection point
$$\vb_r = \bar{\vb} + \rho (\bar{\vb} - \vb_{d + 1}),$$
with $\rho > 0$.
Compute $f(\vb_r)$. \\
Note: Default value for reflection coefficient: $\rho = 1$
\end{enumerate}
\imageC[0.42]{figure/nelder_mead_steps_3.png}
\figcaption{Reflection point $\vb_r$ with $\rho = 1$ (1st iteration)}
\end{framev}


\begin{framev}{nelder-mead method}
Distinguish three cases:\\
\splitVCC[0.57]{
\begin{itemizeS}
\item Case 1: $f(\vb_1) \leq f(\vb_r) < f(\vb_d)$
$\Rightarrow$ Accept $\vb_r$ and discard $\vb_{d + 1}$
\end{itemizeS}
}{\imageC[0.62]{figure/nelder_mead_steps_4.png}
\figcaption{Hypothetical visualization of case 1: acceptance of $\vb_r$}}
\vspace{0.33cm}
\splitVCC[0.57]{
\begin{itemizeS}
\item Case 2: $f(\vb_r) < f(\vb_1)$
$\Rightarrow$ Expansion:
$$\vb_e = \bar{\vb} + \chi (\vb_{r} - \bar{\vb}), \quad \chi > 1$$
We discard $\vb_{d + 1}$ and accept the better of $\vb_r$ and $\vb_e$
(default value for expansion coefficient: $\chi = 2$)
\end{itemizeS}
}{\imageC[0.6]{figure/nelder_mead_steps_5.png}
\figcaption{Visualization of case 2: expansion with $\chi = 2$, but $\vb_r$ is kept since $f(\vb_e) > f(\vb_r)$ (cf. figure above)}}
\end{framev}

\comment[NOTE]{BB}{wir sollten für chi und sigma sagen dass die beide zwischen 0 und 1 sind; nicht 0 bis 0.5 nur für einen von beiden}

\begin{framev}{nelder-mead method}
\begin{itemize}
\item Case 3: $f(\vb_r) \ge f(\vb_d)$
$\Rightarrow$ Contraction:
$$\vb_c = \bar{\vb} + \gamma (\vb_{d + 1} - \bar{\vb})$$
with $0 < \gamma \le 1/2$.\\
Note: default value for contraction coefficient: $\gamma = 1/2$
\begin{itemize}
\item If $f(\vb_c) < f(\vb_{d + 1})$: accept $\vb_c$ and discard $\vb_{d + 1}$
\end{itemize}
\end{itemize}
\splitVThree{
\imageC[0.95]{figure/nelder_mead_steps_6.png}
\figcaption{Reflection (2nd iteration)}
}{
\imageC[0.95]{figure/nelder_mead_steps_7.png}
\figcaption{Case 3: contraction with $\gamma = 1/2$}
}{
\imageC[0.95]{figure/nelder_mead_steps_8.png}
\figcaption{Case 3: acceptance of $\vb_c$}
}
\end{framev}


\begin{framev}{nelder-mead method}
\begin{itemize}
\item Case 3 (continued):
\begin{itemize}
\item If $f(\vb_c) \ge f(\vb_{d + 1})$: shrink entire simplex towards $\vb_1$ (Shrinking)
$$\vb_{i} = \vb_1 + \sigma (\vb_{i} - \vb_{1}) \quad \forall i$$
Note: default value for shrinking coefficient: $\sigma = 1/2$
\end{itemize}
\end{itemize}
\imageC[0.38]{figure/nelder_mead_steps_9.png}
\figcaption{Hypothetical visualization of case 3: shrinkage with $\sigma = 1/2$ (if $f(\vb_c) \ge f(\vb_3)$)}
\begin{enumerate}
\item[4.] Repeat all steps until stopping criterion met
\end{enumerate}
\end{framev}


\begin{framev}{nelder-mead summary}
Advantages:
\begin{itemize}
\item No gradients needed
\item Robust, often works well for non-differentiable functions
\end{itemize}
Drawbacks:
\begin{itemize}
\item Relatively slow (not applicable in high dimensions)
\item Not each step improves, only mean of corner values is reduced
\item No guarantee for convergence to local optimum / stationary point
\end{itemize}
Visualization:
\begin{center}
\url{http://www.benfrederickson.com/numerical-optimization/}
\end{center}
\spacer
Note: Nelder-Mead is default method of \texttt{R} function \texttt{optim()}\\
If gradient is available and cheap, L-BFGS is preferred
\end{framev}


\begin{framev}[align=top]{Nelder-Mead Visualization in 2D}
$$\min_{\xv} f(x_1,x_2) = x_1^{2} + x_2^{2} + x_1\cdot \sin x_2 + x_2 \cdot \sin x_1 $$
\foreach \i in {1, 2, 3, 4}{
\only<\i>{\imageFixed{2.5cm}{2.5cm}[0.9]{figure_man/nm_animation2d_\i.PNG}}
}
\end{framev}


\begin{framev}{Nelder-Mead vs. GD}
\comment[NOTE]{BB}{hier fehlt etwas background info zum bsp: http://www.benfrederickson.com/numerical-optimization/}
\comment[NOTE]{BB}{das bsp könnten wir easy selbt coden. und man muss die objective am besten auch zeigen, die feht nämnlich auf der webseite}
\splitVCC{\imageC[0.85]{figure_man/nm_gd_cities_1.PNG}}{\imageC[0.85]{figure_man/nm_gd_cities_2.PNG}}
\splitVCC{\imageC[0.85]{figure_man/nm_gd_cities_3.PNG}}{\imageC[0.85]{figure_man/nm_gd_cities_4.PNG}}
{\footnotesize Nelder-Mead in multiple dimensions:
Organize points (US cities) to keep predefined mutual distances.
For 10 cities, gradient descent (top) converges well for a suitable learning rate.
Nelder-Mead (bottom) fails to converge, even after many iterations.}
\end{framev}


\begin{framev}{Nelder-Mead vs. GD}
\splitVCC{\imageC[0.85]{figure_man/nm_gd_cities_5.PNG}}{\imageC[0.85]{figure_man/nm_gd_cities_6.PNG}}
\splitVCC{\imageC[0.85]{figure_man/nm_gd_cities_7.PNG}}{\imageC[0.85]{figure_man/nm_gd_cities_8.PNG}}
{\footnotesize Even for only 5 cities, Nelder-Mead (bottom) performs poorly\\
However, gradient descent (top) still works}
\end{framev}

\endlecture
\end{document}
